\section{Background}
\label{sec:background}

\subsection{The Learnable Bridge and the Steersman}

Paper~3 in this series~\cite{bielec2026bridge} introduced the \emph{learnable bridge}:
a per-layer coupling matrix $\mathbf{B}_l \in \mathbb{R}^{n \times n}$ that
mediates interaction between $n$ parallel LoRA channels in a multi-channel
adapter. The bridge is trained jointly with the adapter weights, governed by
a cybernetic feedback loop called the \emph{Steersman}.

The Steersman implements a classical control cycle: a \emph{sensor} extracts
the bridge's graph Laplacian eigenvalues at each feedback interval; an
\emph{oscilloscope} tracks connectivity (Fiedler algebraic connectivity
$\lambda_2$) and co-planar/cross-planar coupling ratios; a \emph{steersman}
module adjusts loss weights; and the modified loss signal acts as the
\emph{actuator} on the bridge parameters. The full loop operates at a
configurable interval (default: every 100 optimizer steps) without
interrupting gradient flow.

Two loss components provide independent control channels:
\begin{itemize}
  \item \textbf{Spectral regularization} targets aggregate connectivity by
    penalizing deviations of $\lambda_2$ from a reference value, ensuring the
    bridge develops non-trivial coupling.
  \item \textbf{Contrastive loss} specifies \emph{which} channels should
    couple preferentially: designated \emph{co-axial} pairs are attracted
    (their coupling is reinforced) while \emph{cross-axial} pairs are repelled
    (their coupling is suppressed).
\end{itemize}

Paper~3 demonstrated that this mechanism, applied with pair specifications
derived from the rhombic dodecahedron's 3-axis face-pair geometry at $n{=}6$,
produces block-diagonal bridge structure with co-planar/cross-planar coupling
ratios exceeding 70{,}000:1. The present paper asks whether this result
generalizes beyond a single polytope.

\subsection{Polytope Geometry and Antipodal Pairing}

A \emph{regular polytope} is a geometric object whose vertices, edges, and
faces are congruent under its symmetry group. Several regular and semi-regular
polytopes admit a natural \emph{antipodal pairing}: for each vertex $v$,
there exists a unique vertex $-v$ at maximal distance through the polytope's
center. When the polytope's vertices are mapped to adapter channels, antipodal
pairs define co-axial specifications for the contrastive loss.

We test four polytope geometries spanning 3D and 4D:

\begin{description}
  \item[Octahedron] ($n{=}4$, 2 axes): The regular octahedron in
    $\mathbb{R}^3$ has 6 vertices forming 3 antipodal pairs. We use 4
    vertices (2 pairs) corresponding to a 2-axis partition.
  \item[Rhombic dodecahedron] ($n{=}6$, 3 axes): The RD's 14 vertices
    include 6 degree-4 vertices forming 3 face-diagonal pairs aligned with
    the coordinate axes.
  \item[Tesseract] ($n{=}8$, 4 axes): The 4D hypercube's 16 vertices form
    8 antipodal pairs. We use 8 vertices (4 pairs) corresponding to the four
    coordinate axes.
  \item[24-cell] ($n{=}24$, 12 axes): The $D_4$ root polytope in
    $\mathbb{R}^4$, self-dual, with 24 vertices as permutations of
    $(\pm 1, \pm 1, 0, 0)$. Its 12 antipodal pairs group into 6
    coordinate-plane families (wx, wy, wz, xy, xz, yz).
\end{description}

The key insight is that the Steersman mechanism requires \emph{no
modification} to accommodate different polytopes. Only the pair specification
--- which channels are co-axial --- changes.

\subsection{Programmable Inductive Biases}

Prior work on structured adapters has explored low-rank
variants~\cite{valipour2023dylora,loraxs2024},
mixture-of-experts~\cite{zadouri2023pushing},
and hypernetworks~\cite{ha2017hypernetworks} that generate adapter weights
conditioned on task embeddings. Zhu et al.~\cite{zhu2024asymmetry} analyze
the asymmetric roles of the $A$ and $B$ projections in LoRA, showing that
structure within the adapter matters beyond rank alone.

Recent work has converged independently on block-diagonal adapter structure.
G{\"u}rses et al.~\cite{gurses2026diablo} prove that block-diagonal
fine-tuning (DiaBlo) is theoretically at least as expressive as standard
LoRA and empirically competitive, validating the block-diagonal form as a
first-class adapter architecture. Wang et al.~\cite{wang2025gralora}
(GraLoRA) and Barazandeh et al.~\cite{barazandeh2025localized} (Localized
LoRA) partition adapter matrices into independent blocks for
parameter-efficient fine-tuning. Tian et al.~\cite{tian2026spectral}
(Spectral Surgery) analyze LoRA's spectral properties for training-free
refinement, demonstrating the diagnostic value of spectral decomposition
in adapter weight spaces.

Our approach differs from all of these in a specific way: we do not merely
\emph{impose} block-diagonal structure (as DiaBlo and GraLoRA do through
architectural constraint) but \emph{program} it through cybernetic
feedback. The pair specification declares which channels are co-axial;
the Steersman compiles this specification into the bridge matrix during
training. The topology is a design choice, not an emergent property or
a fixed architectural constraint. This aligns with geometric deep
learning~\cite{bronstein2021geometric}, where the structure of the
computation graph reflects the structure of the problem domain---but
extends it from the network architecture to the adapter's internal
coupling topology.

\subsection{Spectral Graph Theory Preliminaries}

We characterize bridge structure through the graph Laplacian
$\mathbf{L} = \mathbf{D} - \mathbf{W}$, where $\mathbf{W}$ is the
absolute-value weight matrix derived from the bridge and $\mathbf{D}$ is
the degree matrix. The eigenvalues $0 = \lambda_1 \leq \lambda_2 \leq
\cdots \leq \lambda_n$ encode connectivity:

\begin{itemize}
  \item $\lambda_2$ (Fiedler value): algebraic connectivity. $\lambda_2 = 0$
    indicates a disconnected graph; $\lambda_2 > 0$ indicates a connected
    graph.
  \item Block-diagonal structure manifests as $k$ near-zero eigenvalues
    followed by $k$ large eigenvalues, where $k$ equals the number of
    co-axial pairs --- a clean $k{+}k$ split.
  \item The \emph{co-planar/cross-planar ratio} (co/cross) measures
    directional preference: the ratio of mean absolute coupling within
    co-axial pairs to mean absolute coupling between non-paired channels.
\end{itemize}

These two metrics --- Fiedler value and co/cross ratio --- jointly classify
bridge states into the four regimes described in~\S\ref{sec:negative-controls}.

Tam et al.~\cite{tam2020fiedler} introduced differentiable Fiedler
regularization as a sparsity-promoting objective for neural network
pruning, minimizing $\lambda_2$ to disconnect subgraphs. Our spectral
loss inverts this application: we \emph{target} a specific $\lambda_2$
attractor to program topology, not to prune it.
